% ============================================================
%  专题01  算术平方根与平方根
%  风格：参考 勤思数学 模板
% ============================================================
\documentclass[11pt,a4paper]{ctexart}
\usepackage[margin=2cm,headheight=14pt]{geometry}
\usepackage{amsmath,amssymb}
\usepackage{enumitem}
\usepackage{titlesec}
\usepackage{fancyhdr}
\usepackage{xcolor}
\usepackage{tcolorbox}
\usepackage{booktabs}
\usepackage{needspace}
\usepackage{multicol}
\usepackage{tikz}
\usepackage{pgffor}
\usetikzlibrary{calc}

% === 个人风格配色 ===
\definecolor{primary}{RGB}{0,100,180}      % 主色-深蓝
\definecolor{accent}{RGB}{220,80,60}       % 强调-珊瑚红
\definecolor{teal}{RGB}{0,140,130}         % 辅助-青色
\definecolor{orange}{RGB}{230,130,40}      % 橙色
\definecolor{lightblue}{RGB}{235,245,255}  % 浅蓝背景
\definecolor{lightyellow}{RGB}{255,250,235}
\definecolor{lightgreen}{RGB}{235,250,240}
\definecolor{lightred}{RGB}{255,242,240}

% === 标题样式 ===
\ctexset{
  section = {
    format = \Large\bfseries\color{primary},
    name = {,},
    aftername = \hspace{0.5em},
    beforeskip = 1.2em,
    afterskip = 0.8em,
  },
  subsection = {
    format = \large\bfseries\color{teal},
    beforeskip = 0.8em,
    afterskip = 0.5em,
  }
}

% === 页眉页脚 ===
\pagestyle{fancy}
\fancyhf{}
\fancyhead[L]{\small\color{primary}\bfseries 专题01\ $\cdot$\ 算术平方根与平方根}
\fancyhead[R]{\small\thepage}
\fancyfoot[C]{\small\color{gray}--- 厉飞雨数学 ---}
\renewcommand{\headrulewidth}{0.6pt}
\renewcommand{\footrulewidth}{0.3pt}

% === 自定义盒子 ===
\tcbuselibrary{skins,breakable}

% 知识点盒子
\newtcolorbox{knowledge}[1][]{
  colback=lightblue, colframe=primary,
  fonttitle=\bfseries\color{white},
  title={#1},
  arc=3mm, boxrule=0.8pt,
  left=8pt, right=8pt, top=6pt, bottom=6pt,
  breakable
}

% 例题盒子 —— 白底 + 淡粉标题栏
\newtcolorbox{example}[1][]{
  colback=white, colframe=red!20,
  fonttitle=\bfseries\color{red!50!black},
  colbacktitle=red!10,
  title={#1},
  arc=2mm, boxrule=0.4pt,
  left=8pt, right=8pt, top=5pt, bottom=5pt,
  breakable
}

% 训练盒子 —— 白底 + 淡粉标题栏
\newtcolorbox{practice}[1][]{
  colback=white, colframe=red!20,
  fonttitle=\bfseries\color{red!50!black},
  colbacktitle=red!10,
  title={#1},
  arc=2mm, boxrule=0.4pt,
  left=8pt, right=8pt, top=5pt, bottom=5pt,
  breakable
}

% 易错盒子
\newtcolorbox{warning}[1][]{
  colback=lightred, colframe=accent,
  fonttitle=\bfseries\color{white},
  title={\color{white}#1},
  arc=2mm, boxrule=0.6pt,
  left=8pt, right=8pt, top=5pt, bottom=5pt,
  breakable
}

% === 公式高亮 ===
\newcommand{\key}[1]{\textcolor{accent}{\textbf{#1}}}
\newcommand{\formula}[1]{\textcolor{primary}{\textbf{$#1$}}}

% === 题型编号装饰（robust 防止展开问题）===
\DeclareRobustCommand{\typetag}[1]{%
  \tikz[baseline=(X.base)]{
    \node[fill=primary, text=white, rounded corners=3pt,
          inner xsep=6pt, inner ysep=2pt, font=\small\bfseries] (X) {#1};
  }%
}

% === 答题留白区域 ===

% === 防跨页 ===
\newcommand{\qitem}[1][6]{\needspace{#1\baselineskip}\item}

% === 间距控制 ===
\setlength{\parindent}{0pt}
\setlength{\parskip}{2pt plus 1pt minus 1pt}

% === 选项排列 ===
\setlist[itemize]{nosep,left=0pt .. 2em}

% ============================================================
%  自定义环境
% ============================================================

% ── 选择题选项 —— 两列清爽排版 ────────────────────────────
\newcommand{\choices}[4]{%
  \par\vspace{3pt}%
  \noindent
  \begin{minipage}[t]{0.47\textwidth}A．#1\end{minipage}%
  \hfill
  \begin{minipage}[t]{0.47\textwidth}B．#2\end{minipage}\\[4pt]
  \noindent
  \begin{minipage}[t]{0.47\textwidth}C．#3\end{minipage}%
  \hfill
  \begin{minipage}[t]{0.47\textwidth}D．#4\end{minipage}%
  \par\vspace{2pt}%
}

% ── 子题网格 —— 2×2 对齐排列 ──────────────────────────────
\newcommand{\subqs}[4]{%
  \par\vspace{3pt}%
  \noindent
  \begin{minipage}[t]{0.47\textwidth}#1\end{minipage}%
  \hfill
  \begin{minipage}[t]{0.47\textwidth}#2\end{minipage}\\[5pt]
  \noindent
  \begin{minipage}[t]{0.47\textwidth}#3\end{minipage}%
  \hfill
  \begin{minipage}[t]{0.47\textwidth}#4\end{minipage}%
  \par\vspace{2pt}%
}
\newcommand{\subqsThree}[3]{%
  \par\vspace{3pt}%
  \noindent
  \begin{minipage}[t]{0.47\textwidth}#1\end{minipage}%
  \hfill
  \begin{minipage}[t]{0.47\textwidth}#2\end{minipage}\\[5pt]
  \noindent
  \begin{minipage}[t]{0.47\textwidth}#3\end{minipage}%
  \par\vspace{2pt}%
}
\newcommand{\subqsFive}[5]{%
  \par\vspace{3pt}%
  \noindent
  \begin{minipage}[t]{0.47\textwidth}#1\end{minipage}%
  \hfill
  \begin{minipage}[t]{0.47\textwidth}#2\end{minipage}\\[5pt]
  \noindent
  \begin{minipage}[t]{0.47\textwidth}#3\end{minipage}%
  \hfill
  \begin{minipage}[t]{0.47\textwidth}#4\end{minipage}\\[5pt]
  \noindent
  \begin{minipage}[t]{0.47\textwidth}#5\end{minipage}%
  \par\vspace{2pt}%
}

% ── 填空留白 ───────────────────────────────────────────────
\newcommand{\fillblank}[1][2cm]{\underline{\hspace{#1}}}

% ── 解答留白 —— 纯空白行 ──────────────────────────────────
% 用法: \solutionblank（默认2行）或 \solutionblank[4]（4行）
\newcommand{\solutionblank}[1][2]{%
  \par\vspace{4pt}%
  \foreach \i in {1,...,#1}{%
    \par\vspace{10pt}%
  }%
  \vspace{-2pt}%
}

% ============================================================
\begin{document}

% === 封面 ===
\begin{titlepage}
\begin{center}
\vspace*{3cm}

{\tikz{
  \node[fill=primary, text=white, rounded corners=8pt,
        inner xsep=20pt, inner ysep=10pt,
        font=\Huge\bfseries] {算术平方根与平方根};
}}

\vspace{1.5cm}

{\Large\color{gray} 专题01\ $\cdot$\ 七年级数学}

\vspace{2cm}

\begin{tikzpicture}
  \node[draw=primary, line width=1pt, rounded corners=5pt,
        inner sep=15pt, text width=10cm, align=left,
        font=\large] {
    \textcolor{primary}{\textbf{本讲要点}}\\[8pt]
    \textcolor{teal}{$\bullet$} 算术平方根的定义与性质\\[4pt]
    \textcolor{teal}{$\bullet$} 平方根的概念与性质\\[4pt]
    \textcolor{teal}{$\bullet$} 算术平方根与平方根的区别与联系\\[4pt]
    \textcolor{teal}{$\bullet$} 估算与实际应用
  };
\end{tikzpicture}

\vfill
{\large\color{gray} \today}
\end{center}
\end{titlepage}

% ========================================
% 知识点部分
% ========================================

\section*{\color{primary}\Large $\blacktriangleright$\ 知识要点}

% --- 知识点1 ---
\begin{knowledge}[\color{white}知识点1\ $\cdot$\ 算术平方根]
\textbf{定义}：一般地，如果一个正数 $x$ 的平方等于 $a$，即 $\formula{x^2=a}$，那么这个正数 $x$ 叫作 $a$ 的\key{算术平方根}。

\vspace{0.3em}
\textbf{表示方法}：$a$ 的算术平方根记为 $\formula{\sqrt{a}}$，读作"根号 $a$"，$a$ 叫做\textbf{被开方数}。
实际上 $\sqrt{\phantom{a}}$ 中省略了根指数 $2$，因此也读作"二次根号 $a$"。

\vspace{0.3em}
\textbf{性质}：
\begin{enumerate}[leftmargin=1.8em]
  \item 算术平方根是它本身的数有 \key{$0$ 和 $1$}；
  \item 一个正数的算术平方根\textbf{只有一个}，并且\textbf{恒为正}；$0$ 的算术平方根为 $0$，即 $\sqrt{0}=0$；\textbf{负数没有算术平方根}；
  \item 具有\key{双重非负性}：被开方数 $\sqrt{a}\geq 0$（$\geq 0$），其本身 $\sqrt{a}\geq 0$（$\geq 0$）。
\end{enumerate}
\end{knowledge}

% --- 知识点2 ---
\begin{knowledge}[\color{white}知识点2\ $\cdot$\ 小数点移动规律]
一个数扩大为原来的 $100$ 倍，它的算术平方根就扩大为原来的 $10$ 倍。
一个数缩小为原来的 $\dfrac{1}{100}$，它的算术平方根就缩小为原来的 $\dfrac{1}{10}$。

\begin{center}
\begin{tikzpicture}
  \node[draw=primary, fill=lightblue, rounded corners=5pt,
        inner sep=8pt, text width=12cm, align=center, font=\small] {
    \textbf{规律}：被开方数的小数点向右（或向左）移动\textbf{两位}，
    它的算术平方根的小数点相应地向右（或向左）移动\textbf{一位}。
  };
\end{tikzpicture}
\end{center}
\end{knowledge}

% --- 知识点3 ---
\begin{knowledge}[\color{white}知识点3\ $\cdot$\ 平方根]
\textbf{定义}：一般地，如果一个数 $x$ 的平方等于 $a$，即 $\formula{x^2=a}$，那么这个数 $x$ 叫作 $a$ 的\key{平方根}（也称为\textbf{二次方根}）。

\vspace{0.3em}
\textbf{表示方法}：正数 $a\,(a\geq 0)$ 的平方根记作 $\formula{\pm\sqrt{a}}$，读作"正负根号 $a$"。
其中 $\sqrt{a}$ 表示正平方根（算术平方根），$-\sqrt{a}$ 表示负平方根。$0$ 的平方根记作 $\sqrt{0}=0$。

\vspace{0.3em}
\textbf{性质}：
\begin{itemize}[leftmargin=1.8em]
  \item 一个\textbf{正数}有两个平方根，它们\key{互为相反数}；
  \item $0$ 只有一个平方根，是它本身；
  \item \textbf{负数}没有平方根。
\end{itemize}

\vspace{0.3em}
\textbf{开平方}：求一个数 $a$ 的平方根的运算叫作\key{开平方}。
开平方与平方\textbf{互为逆运算}。
\end{knowledge}

% --- 知识点4 ---
\begin{knowledge}[\color{white}知识点4\ $\cdot$\ 算术平方根 vs 平方根]
\begin{center}
\begin{tabular}{lcc}
  \toprule
  & \textbf{算术平方根} & \textbf{平方根} \\
  \midrule
  表示 & $\sqrt{a}$ & $\pm\sqrt{a}$ \\
  个数 & 正数有 \textbf{1} 个 & 正数有 \textbf{2} 个 \\
  正负 & $\geq 0$ & 一正一负（$0$ 除外） \\
  被开方数要求 & $a\geq 0$ & $a\geq 0$ \\
  关系 & 正数的正平方根 $=$ 算术平方根 & 平方根 $=\pm$ 算术平方根 \\
  \bottomrule
\end{tabular}
\end{center}
\end{knowledge}

\begin{warning}[\color{white}易错警示]
\begin{itemize}[leftmargin=1.5em]
  \item $\sqrt{a}$ 表示的是\textbf{算术平方根}，结果恒为非负数。
  \item "$\pm\sqrt{a}$" 才表示平方根，包含正负两个值。
  \item $a$ 的算术平方根 $\sqrt{a}$ 与 $a$ 的平方根 $\pm\sqrt{a}$ 是\textbf{不同的概念}。
  \item $0$ 的平方根 $=$ 算术平方根 $=$ $0$。
\end{itemize}
\end{warning}

\newpage

% ============================================================
\subsection*{\typetag{题型01}\ 辨析算术平方根的概念}
% ============================================================

\begin{example}
\textbf{【例1】} 下列说法正确的是（\quad\quad）

\choices{$\sqrt{25}$ 表示 $25$ 的算术平方根}{$\sqrt{2}$ 表示 $2$ 的算术平方根}{$2$ 的算术平方根记作 $\sqrt{2}$}{$2$ 是 $\sqrt{4}$ 的算术平方根}
\end{example}

\begin{practice}[\color{teal}跟踪训练]
\textbf{1.} 算术平方根是它本身的数是（\quad\quad）

\medskip
\textbf{2.} 如果 $\sqrt{a}$ 有算术平方根，那么 $a$ 可以取的值为（\quad\quad）

\medskip
\textbf{3.} 下列各数没有算术平方根的是（\quad\quad）

\end{practice}

% ============================================================
\subsection*{\typetag{题型02}\ 求算术平方根}
% ============================================================

\begin{example}
\textbf{【例2】} $\sqrt{4}$ 的算术平方根是（\quad\quad）

\medskip
\textbf{【例3】} 计算：$\sqrt{\dfrac{4}{9}}=\underline{\hspace{1.5cm}}$，$\sqrt{0.01}=\underline{\hspace{1.5cm}}$。

\medskip
\textbf{【例4】} $\sqrt{81}$ 的算术平方根是 $\underline{\hspace{1.5cm}}$。
\end{example}

\begin{practice}[\color{teal}跟踪训练]

\medskip

\medskip

\medskip
\end{practice}

% ============================================================
\subsection*{\typetag{题型03}\ 辨析平方根的概念}
% ============================================================

\begin{example}
\textbf{【例5】} "$\dfrac{4}{9}$ 的平方根是 $\pm\dfrac{2}{3}$"，用数学式子表达为（\quad\quad）

\choices{$\sqrt{\dfrac{4}{9}}=\pm\dfrac{2}{3}$}{$\pm\sqrt{\dfrac{4}{9}}=\pm\dfrac{2}{3}$}{$\pm\sqrt{\dfrac{4}{9}}=\dfrac{2}{3}$}{$\sqrt{\dfrac{4}{9}}=\dfrac{2}{3}$}

\medskip
\textbf{【例6】} 下列各数中没有平方根的是（\quad\quad）

\medskip
\textbf{【例7】} 下列说法正确的是（\quad\quad）

\end{example}

\begin{practice}[\color{teal}跟踪训练]
\textbf{1.} 下列语句写成数字式子正确的是（\quad\quad）

\choices{$9$ 是 $81$ 的算术平方根：$9=\sqrt{81}$}{$5$ 是 $25$ 的算术平方根：$5=\pm\sqrt{25}$}{$-6$ 是 $36$ 的平方根：$-6=\sqrt{36}$}{$-2$ 是 $4$ 的负的平方根：$-2=-\sqrt{4}$}

\medskip
\textbf{2.} 下列说法正确的是（\quad\quad）

\medskip
\textbf{3.} 下列说法正确的是（\quad\quad）

\medskip
\textbf{4.} 一个自然数的一个平方根是 $a$，则与它相邻的下一个自然数的平方根是（\quad\quad）

\end{practice}

% ============================================================
\subsection*{\typetag{题型04}\ 求平方根}
% ============================================================

\begin{example}
\textbf{【例8】} 求下列各数的平方根：

\subqs{(1)\,$2.25$}{(2)\,$\dfrac{49}{64}$}{(3)\,$0$}{(4)\,$-16$}

\medskip

\medskip
\textbf{【例10】} $\sqrt{81}$ 的平方根是（\quad\quad）

\end{example}

\begin{practice}[\color{teal}跟踪训练]

\medskip
\textbf{2.} 求下列各数的平方根，并用式子表示出来：

\subqs{(1)\,$100$}{(2)\,$0.49$}{(3)\,$\dfrac{9}{16}$}{(4)\,$(-5)^{2}$}

\medskip
\textbf{3.} 求下列等式中的 $x$：

\subqs{(1)\,若 $x^{2}=49$，则 $x=\underline{\hspace{1cm}}$}{(2)\,若 $x^{2}=\dfrac{1}{4}$，则 $x=\underline{\hspace{1cm}}$}{(3)\,若 $x^{2}=(-3)^{2}$，则 $x=\underline{\hspace{1cm}}$}{(4)\,若 $x^{2}=0$，则 $x=\underline{\hspace{1cm}}$}
\end{practice}

% ============================================================
\subsection*{\typetag{题型05}\ 利用概念求值}
% ============================================================

\begin{example}
\textbf{【例11】} 已知 $2a-1$ 的平方根为 $\pm 3$，$3a+b$ 的算术平方根为 $4$。

\hspace{2em}(1) 求 $a$、$b$ 的值；\quad (2) 求 $a+b$ 的平方根。
\solutionblank[4]
\end{example}

\begin{practice}[\color{teal}跟踪训练]
\textbf{1.} 已知 $2a-b^{2}$ 的平方根是 $\pm 4$，$a-2b$ 的平方根是 $\pm 2$，求 $a+b$ 的平方根。
\solutionblank[4]

\medskip
\textbf{2.} 已知 $a$ 的值是 $2$，$b$ 的算术平方根是 $4$。

\hspace{2em}(1) 求 $a$、$b$ 的值；\quad (2) 求 $a+b$ 的平方根。
\solutionblank[4]

\medskip
\textbf{3.} 已知 $x=3$，$y=4$，$z$ 是 $9$ 的算术平方根，求 $x+y-z$ 的平方根。
\solutionblank[3]
\end{practice}

% ============================================================
\subsection*{\typetag{题型06}\ 利用平方根解方程}
% ============================================================

\begin{example}
\textbf{【例12】} 求下列各式中的值：

\hspace{2em}(1)\,$(x-1)^{2}=25$；\quad (2)\,$(2x+1)^{2}=9$；\quad (3)\,$4x^{2}-49=0$。
\solutionblank[4]
\end{example}

\begin{practice}[\color{teal}跟踪训练]
\textbf{1.} 求下列各式中 $x$ 的值：

\subqsThree{(1)\,$x^{2}=81$}{(2)\,$(x-2)^{2}=16$}{(3)\,$(3x+1)^{2}=25$}
\solutionblank[5]

\medskip
\textbf{2.} 解方程：

\subqsThree{(1)\,$(x-3)^{2}=49$}{(2)\,$(2x-1)^{2}=16$}{(3)\,$9x^{2}-1=0$}
\solutionblank[5]
\end{practice}

% ============================================================
\subsection*{\typetag{题型07}\ 非负性的应用}
% ============================================================

\begin{example}

\medskip
\textbf{【例14】} 已知 $a$、$b$ 都是实数，且 $\sqrt{a-2}+|b+3|=0$，求 $a+b$ 的平方根。
\solutionblank[3]
\end{example}

\begin{practice}[\color{teal}跟踪训练]

\medskip
\end{practice}

% ============================================================
\subsection*{\typetag{题型08}\ 平方根性质的应用}
% ============================================================

\begin{example}
\end{example}

\begin{practice}[\color{teal}跟踪训练]

\medskip
\end{practice}

% ============================================================
\subsection*{\typetag{题型09}\ 运算}
% ============================================================

\begin{example}
\textbf{【例16】} 计算：

\subqsFive{(1)\,$\sqrt{64}$}{(2)\,$\sqrt{(-5)^{2}}$}{(3)\,$-\sqrt{0.09}$}{(4)\,$\sqrt{(-13)^{2}}$}{(5)\,$\sqrt{16}+\sqrt{9}$}

\medskip
\textbf{【例17】} 求下列各式的值：

\subqs{(1)\,$\sqrt{(-8)^{2}}$}{(2)\,$\sqrt{(-0.5)^{2}}$}{(3)\,$\sqrt{\left(\dfrac{3}{7}\right)^{2}}$}{(4)\,$\pm\sqrt{\dfrac{25}{36}}$}

\medskip
\textbf{【例18】} 求下列各式的值：

\subqsThree{(1)\,$\sqrt{49}-\sqrt{16}$}{(2)\,$\sqrt{\dfrac{9}{4}}+\sqrt{\dfrac{1}{9}}$}{(3)\,$\sqrt{1}+\sqrt{4}+\sqrt{9}$}
\end{example}

\begin{practice}[\color{teal}跟踪训练]
\textbf{1.} 下列选项中，正确的是（\quad\quad）

\choices{$\sqrt{(-3)^{2}}=-3$}{$\sqrt{16}=\pm 4$}{$-\sqrt{25}=-5$}{$\sqrt{0}=1$}

\medskip
\textbf{2.} 求下列各式的值：

\subqs{(1)\,$\sqrt{25}$}{(2)\,$\sqrt{0.01}$}{(3)\,$\sqrt{\dfrac{49}{81}}$}{(4)\,$\pm\sqrt{100}$}

\medskip
\textbf{3.} 求下列各式的值：

\noindent(1)\,$\sqrt{(-9)^{2}}=\underline{\hspace{1cm}}$\quad (2)\,$\sqrt{(-4)^{2}}=\underline{\hspace{1cm}}$\quad (3)\,$\pm\sqrt{\dfrac{4}{25}}=\underline{\hspace{1cm}}$\quad (4)\,$\sqrt{0.16}=\underline{\hspace{1cm}}$

\vspace{4pt}
\noindent(5)\,$\sqrt{1.44}=\underline{\hspace{1cm}}$\quad (6)\,$-\sqrt{36}=\underline{\hspace{1cm}}$\quad (7)\,$\pm\sqrt{\dfrac{9}{4}}=\underline{\hspace{1cm}}$\quad (8)\,$\sqrt{625}=\underline{\hspace{1cm}}$
\end{practice}

% ============================================================
\subsection*{\typetag{题型10}\ 估算}
% ============================================================

\begin{example}
\textbf{【例19】} 若 $\sqrt{a}$ 的值在 $3$ 与 $4$ 之间，则满足条件的 $a$ 可能是（\quad\quad）

\medskip
\end{example}

\begin{practice}[\color{teal}跟踪训练]
\textbf{1.} 估计 $\sqrt{20}$ 的值在（\quad\quad）

\medskip
\textbf{2.} 已知 $a$，$b$ 是两个连续整数，$a<\sqrt{7}<b$，则 $a+b$ 的值为（\quad\quad）

\medskip
\textbf{3.} 一个正方形的面积是 $8$，估计它的边长大小在（\quad\quad）

\medskip

\medskip
\end{practice}

% ============================================================
\subsection*{\typetag{题型11}\ 小数点移动规律}
% ============================================================

\begin{example}
\textbf{【例21】} 已知 $\sqrt{3}\approx 1.732$，那么下列结论正确的是（\quad\quad）

\medskip

\begin{center}
\begin{tabular}{c|cccccc}
\toprule
$a$ & $\cdots$ & $0.000001$ & $0.01$ & $1$ & $100$ & $10000$ \\
\midrule
$\sqrt{a}$ & $\cdots$ & $0.001$ & $0.1$ & $1$ & $10$ & $100$ \\
\bottomrule
\end{tabular}
\end{center}
\end{example}

\begin{practice}[\color{teal}跟踪训练]
\textbf{1.} 若 $\sqrt{3}\approx 1.732$，$\sqrt{30}\approx 5.477$，则 $\sqrt{300}$ 的平方根约为（\quad\quad）

\medskip
\textbf{2.} (1) 填表并观察规律：

\begin{center}
\begin{tabular}{c|cccccc}
\toprule
$a$ & $0.0064$ & $0.64$ & $64$ & $6400$ \\
\midrule
$\sqrt{a}$ & \rule{2cm}{0.4pt} & \rule{2cm}{0.4pt} & \rule{2cm}{0.4pt} & \rule{2cm}{0.4pt} \\
\bottomrule
\end{tabular}
\end{center}

(2) 根据你发现的规律填空：

\hspace{2em}① 已知 $\sqrt{0.5}\approx 0.707$，则 $\sqrt{50}\approx\underline{\hspace{2cm}}$；

\hspace{2em}② 已知 $\sqrt{5}\approx 2.236$，则 $\sqrt{0.005}\approx\underline{\hspace{2cm}}$。

(3) 从以上问题的解决过程中，你发现了什么规律，试简要说明。
\solutionblank[3]

\medskip
\textbf{3.} 按要求填空：

(1) 填表并观察规律：

\begin{center}
\begin{tabular}{c|ccccc}
\toprule
$a$ & & & $4$ & $400$ \\
\midrule
$\sqrt{a}$ & & & & \\
\bottomrule
\end{tabular}
\end{center}

(2) 根据你发现的规律填空：

\hspace{2em}已知：$\sqrt{0.4}\approx 0.632$，则 $\sqrt{40}\approx\underline{\hspace{2cm}}$；

\hspace{2em}已知：$\sqrt{4}\approx 2$，$\sqrt{400}=20$，则 $\sqrt{40000}\approx\underline{\hspace{2cm}}$。

(3) 从以上问题的解决过程中，你发现了什么规律，试简要说明。
\solutionblank[3]
\end{practice}

% ============================================================
\subsection*{\typetag{题型12}\ 实际应用}
% ============================================================

\begin{example}
\textbf{【例23】} 如图，用两个边长为 $a$ 的小正方形拼成一个大的正方形。


\hspace{2em}(2) 若沿此大正方形边长的方向剪出一个长方形，能否使剪出的长方形纸片的长宽之比为 $3:2$ 且面积为 $24\,\text{cm}^{2}$？
\solutionblank[4]
\end{example}

\begin{practice}[\color{teal}跟踪训练]
\textbf{1.} 观察正方形方格，每个小正方形的边长均是 $1$。

(1) 如图1，求阴影正方形的面积和边长；

(2) 图2是 $10\times 10$ 的正方形方格，请在图2中画出长为 $\sqrt{5}$ 的线段，并说明理由。
\solutionblank[5]

\medskip
\textbf{2.} 为了装饰房间，小明制作了一个面积为 $4\,\text{m}^{2}$ 的正方形拼图。他准备把这个拼图装进一个长方形相框中，这个长方形相框的长和宽之比为 $3:1$，且面积为 $8\,\text{m}^{2}$。

(1) 求长方形相框的长和宽。

(2) 小明能将拼图放入这个相框中吗？请通过计算说明。
\solutionblank[5]

\medskip
\textbf{3.} 如图，用两个面积为 $25\,\text{cm}^{2}$ 的小正方形剪拼成一个大的正方形。


(2) 若沿此大正方形纸片的边的方向剪出一个长方形，能否使剪出的长方形纸片的长、宽之比为 $3:2$，且面积为 $24\,\text{cm}^{2}$？若能，试求出剪出的长方形纸片的长宽；若不能，试说明理由。

（参考数据：$\sqrt{2}\approx 1.414$，$\sqrt{3}\approx 1.732$）
\solutionblank[5]
\end{practice}

\newpage

% ============================================================
%  专题练习
% ============================================================
\begin{center}
  \vspace{2pt}
  {\LARGE\sffamily\bfseries 专题练习}\\[4pt]
  {\color{blue!30}\rule{\textwidth}{0.6pt}}
\end{center}
\vspace{8pt}

% ── 一、选择题 ──────────────────────────────────────────────

\textbf{1.}（2023 上海闵行区期中）$\sqrt{16}$ 的算术平方根等于（\quad\quad）

\medskip
\textbf{2.}（2022 上海静安 期中）$9$ 的平方根为（\quad\quad）

\medskip
\textbf{3.}（2024 静安区七年级下期中）下列说法正确的是

\choices{一个数的平方根一定有两个}{一个正数的平方根一定是它的算术平方根}{一个非零数的正的平方根一定是它的算术平方根}{一个非负数的非负平方根一定是它的算术平方根}

\medskip
\textbf{4.}（2024 松江区七年级下期中）面积为 $9$ 的正方形，其边长等于（\quad\quad）

\medskip
\textbf{5.}（2024 上海实验七年级期中）下列说法正确的是（\quad\quad）

\choices{$4$ 是 $16$ 的一个平方根}{$16$ 的平方根是 $4$}{$4$ 的算术平方根是 $\pm 2$}{$16$ 的算术平方根是 $4$}

\medskip
\textbf{6.}（2024 存志中学期中）下列说法正确的是（\quad\quad）

\choices{平方根等于它本身的数是 $0$，$1$}{倒数等于它本身的数只有 $1$}{算术平方根等于它本身的数是 $0$，$1$}{$4$ 的平方根为 $2$}

% ── 二、填空题 ──────────────────────────────────────────────


\medskip

\medskip

\medskip

\medskip

\medskip

\medskip

\medskip
\textbf{14.}（2024 静安区七年级下期中）若 $\sqrt{x-3}+|y+2|=0$，则 $\underline{\hspace{1cm}}$。

\medskip

\medskip

\medskip
\textbf{17.}（2024 浦东新区七年级下期中）已知 $(x-1)^{2}=4$，则 $\underline{\hspace{1cm}}$。

\medskip

% ── 三、解答题 ──────────────────────────────────────────────

\begin{example}
\textbf{19.}（2024 徐汇中学七年级练习）求下列各数的算术平方根：

\subqs{(1)\,$169$}{(2)\,$\dfrac{49}{64}$}{(3)\,$0.09$}{(4)\,$(-5)^{2}$}
\solutionblank[4]
\end{example}

\begin{example}
\textbf{20.} 求下列各数的平方根：

\subqs{(1)\,$121$}{(2)\,$0.01$}{(3)\,$\dfrac{16}{25}$}{(4)\,$(-7)^{2}$}
\solutionblank[4]
\end{example}

\begin{example}
\textbf{21.} 求下列各式中的 $x$ 值：

\subqsFive{(1)\,$x^{2}=64$}{(2)\,$(x-1)^{2}=9$}{(3)\,$(2x+3)^{2}=25$}{(4)\,$4(x-2)^{2}=36$}{(5)\,$(3x-1)^{2}=16$}
\solutionblank[5]
\end{example}

\begin{example}
\textbf{22.}（21-22 八年级 全国 假期作业）已知 $\sqrt{x-y+3}$ 与 $(2x+3y+1)^{2}$ 互为相反数，求 $x-y$ 的平方根。
\solutionblank[4]
\end{example}

\begin{example}
\textbf{23.}（2023 青浦区七年级统考期末）已知实数 $a$ 的平方根为 $3a+2$ 和 $a-6$，$b$ 的整数部分为 $b$。

\hspace{2em}(1) 求 $a$，$b$ 的值；

\hspace{2em}(2) 若 $b$ 的小数部分为 $c$，求 $a+b+c$ 的平方根。
\solutionblank[5]
\end{example}

\begin{example}
\textbf{24.}（24-25 七年级下 宝山期中）已知 $x=2$，$y=3$，$z$ 是 $9$ 的算术平方根，求 $x+y-z$ 的值。
\solutionblank[3]
\end{example}

\begin{example}
\textbf{25.}（24-25 七年级下 上外附中阶段练习）已知 $2m-4$ 的平方根是 $\pm 2$，$3m+n$ 的平方根是 $\pm 4$。

\hspace{2em}(1) 求 $m$，$n$ 的值；

\hspace{2em}(2) 求 $m+n$ 的平方根。
\solutionblank[4]
\end{example}

\begin{example}
\textbf{26.}（24-25 七年级下 上宝中学阶段练习）一个正数的两个不同的平方根分别是 $3a+2$ 和 $-a+4$。

\hspace{2em}(1) 求 $a$ 和这个正数的值；

\hspace{2em}(2) 求这个数的算术平方根。
\solutionblank[4]
\end{example}

\begin{example}
\textbf{27.}（23-24 七年级下 华育中学月考）如图，用两个面积为 $25\,\text{cm}^{2}$ 的小正方形纸片拼成一个大正方形。

\hspace{2em}(1) 求拼成的大正方形纸片的边长；

\hspace{2em}(2) 小丽想：若沿此大正方形纸片的边的方向剪出一个长方形，能否使剪出的长方形纸片的长、宽之比为 $3:2$ 且面积为 $24\,\text{cm}^{2}$？她不知能否剪得出来，正在发愁。小明见了说："别发愁，一定能用一块面积大的纸片剪出一块面积小的纸片。"你同意小明的说法吗？你认为小丽能用这块纸片剪出符合要求的纸片吗？为什么？
\solutionblank[6]
\end{example}

\end{document}
